
dont ils sont déduits est compatible avec les structures d'anneaux, et grâce à (7.5.1), on trouve que l'on a
\[
  \langle\sigma_\bullet(x_t),\sigma^\bullet(y_t)\rangle
  =
  \langle x_t,y_t\rangle
  \qquad\text{pour }
  x_t\in\Gr_\bullet(X_t)\quad y_t\in\Gr^\bullet(X_t),
\tag{7.9.3}
\]
en utilisant les accouplements \(\langle x,y\rangle\) de 4.1 entre \(\Gr_\bullet\) et \(\Gr^\bullet\). De ceci on déduit formellement que \emph{si} \(x\) \emph{est un élément de} \(\Gr_\bullet(X_t)\) (\emph{resp}. \emph{de} \(\Gr^\bullet(X_t)\)) \emph{tel que} \(\sigma_\bullet(x)\) (\emph{resp}. \(\sigma^\bullet(x)\)) \emph{soit numériquement équivalent à zéro} (4.5.1), \emph{alors il en est de même de} \(x\). La réciproque est vraie si (7.9.2) (ou ce qui revient au même, l'un ou l'autre homomorphisme (7.9.1) qui s'en déduit par passage aux gradués associés) est un épimorphisme mod.torsion, - circonstance qui n'est cependant réalisée qu'assez exceptionnellement.

\textbf{7.9.4.} Il est plausible, du moins lorsque \(f\) est propre et \emph{lisse}, que la réciproque envisagée soit toujours valable, i.e. que \(\sigma_\bullet(x)\) (resp. \(\sigma^\bullet(x)\)) soit numériquement équivalent à zéro, dès qu'il en est ainsi de \(x\), ce qui impliquerait que les homomorphismes (7.9.1) induisent des \emph{monomorphismes} pour les groupes quotients correspondants par la relation d'équivalence numérique \(\Gr^{\mathrm{num}}_\bullet\) et \(\Gr_{\mathrm{num}}^\bullet\). Du moins ce qui précède implique que le groupe \(\Gr_\bullet^{\mathrm{num}}(X_t)\) est isomorphe à un quotient d'un sous-groupe de \(\Gr_\bullet^{\mathrm{num}}(X_s)\) (et en particulier, il a un rang au plus égal à celui de \(\Gr_\bullet^{\mathrm{num}}(X_s)\)), et de même pour \(\Gr_{\mathrm{num}}^\bullet\), et aussi pour \(K^\bullet_{\mathrm{num}}\) lui-même.

\subsection*{7.10. Compatibilité avec les images directes et images inverses}
Donnons-nous un morphisme
\[
  g:X\longrightarrow Y
\]
de schémas noethériens et plats sur \(S\). On désigne par \(\sigma_X\) (resp. \(\sigma_Y\)) l'homomorphisme de spécialisation pour \(X\) (resp. \(Y\)), et par \(i_X,j_X\) (resp. \(i_Y,j_Y\))
